\documentclass[10pt,a4paper]{amsart}
\usepackage[margin=19mm]{geometry}
\usepackage{amsmath,amssymb,mathtools}
\usepackage[scr=rsfs,cal=euler]{mathalfa}
\usepackage{xfrac}
\usepackage[unicode,hidelinks,bookmarks=false]{hyperref}
\newcommand{\ie}{i.e.\ }
\newcommand{\cf}{cf.\ }
\newcommand{\resp}{respectively\ }
\newcommand{\Subsection}{\subsection}
\providecommand{\nobreakdash}{\nobreak\hbox{-}}
\setlength{\emergencystretch}{2em}
\multlinegap0pt
\usepackage{xcolor}
\newtheorem{thm}{Theorem}
\title{The versal deformation of elliptic \textit{m}-fold point curve singularities}
\author{Jan Stevens}
\date{Source: arXiv:2504.16569v1 (CC BY 4.0)}
\begin{document}
\maketitle

\noindent\emph{Source and licence.} The versal deformation of elliptic \textit{m}-fold point curve singularities. Version arXiv:2504.16569v1 (CC BY 4.0), distributed by arXiv under the Creative Commons Attribution 4.0 International licence, http://creativecommons.org/licenses/by/4.0/, as recorded in the arXiv licence metadata for this exact version and shown on the abstract page https://arxiv.org/abs/2504.16569v1. Full licence notice: \texttt{licenses/arxiv-2504.16569-CC-BY-4.0.txt}.\par\smallskip

\noindent\textit{Editorial scope.} Counted unit: the article's thm thm:basetotal (source0.tex lines 711--716). The article's complete original proof of it is delivered with it (source0.tex lines 717--738, one \texttt{proof} environment of 22 lines). No mathematics is written, repaired or re-derived: the statement and the proof are the article's own lines, in the article's own notation, and no retained line is substituted or re-flowed.  Retained uncounted as the source-local context the proof needs: lines 619--635; lines 702--704.  Nothing was omitted from any retained span, so the package carries no omission record.  No retained line contains a citation, so the wrapper carries no bibliography and no bibliography entry.  No retained line contains a figure environment, an image-inclusion command or drawing code, so no image is delivered and the package declares no asset dependency.  Editorial formatting only, in addition: an AMS \texttt{amsart} preamble that restates the article's own declarations and macros used by the retained lines, namely the environments \texttt{thm} on separate counters, together with the standard packages those lines need.  The retained environments print in this document as Theorem 1 in order of appearance, whereas the article prints the same items with the numbering of its own full document.  The complete native source of the article as distributed in the arXiv e-print for this exact version is bundled unchanged as \texttt{sources/176961/source0.tex}.
\begin{thm}\label{versal-def-thm}
The versal deformation of $L_{n+1}^n$, where $n\geq 4$, is given by the polynomials
\begin{align*}
%\quad 
F_{ij;il}&= z_iz_j-z_iz_l+(a_{ij}-a_{il})z_i+a_{ji}z_j-a_{li}z_l \\
&\qquad - a_{ij}a_{ik}-a_{ji}a_{jk}-a_{ki}a_{kj}+ a_{il}a_{ik}+a_{li}a_{lk}+a_{ki}a_{kl}\\
&=(z_i-a_{ij})(z_j-a_{ji})-(a_{ik}-a_{ij})(a_{jk}-a_{ji})\\
&\qquad -(z_i-a_{il})(z_l-a_{li})+(a_{ik}-a_{il})(a_{lk}-a_{li})\;,
\end{align*}
where $a_{ij}+a_{ji}=0$ for all $1\leq i,j\leq n$,
with base space $B_n$ given by 
\begin{multline*}
\quad \Phi^i_{jl;km}=(a_{ik}-a_{ij})(a_{jk}-a_{ji})-(a_{ik}-a_{il})(a_{lk}-a_{li})\\
- (a_{im}-a_{ij})(a_{jm}-a_{ji})+(a_{im}-a_{il})(a_{lm}-a_{li})\quad
\end{multline*}
for all distinct $i$, $j$, $k$, $l$, $m$.
\end{thm}  

\begin{equation}\label{phi-eqn}
\Phi^i_{jl;km}-\Phi^n_{jl;km}-\Phi^k_{jl;in}+\Phi^m_{jl;in}=0\;.
\end{equation}

\begin{thm}\label{thm:basetotal}
For $n\geq 5$ 
the base space $B_n$ of the versal deformation of $L_{n+1}^n$
is  isomorphic to 
the total space of the versal deformation of $L_{n}^{n-1}$.
\end{thm}

\begin{proof}
We start from the equations of the total space and the base space
in Theorem \ref{versal-def-thm}.
By substituting $z_i=a_{in}$ the polynomial $F_{ij;il}$ (written with the index $k$)
%from Theorem \ref{versal-def-thm} 
becomes  the polynomial $\Phi^i_{jl;nk}$. 
We establish by induction that this procedure gives (together with the equations for the base space of $L_{n}^{n-1}$) equations describing the base space of $L_{n+1}^n$,
by showing that the number of independent equations  is equal to the
dimension of $T^2$, which is $\frac{n(n+1)(n-4)}6$.

The polynomials $F_{ij;il}$  give
$\frac{n(n-3)}2$ linearly independent equations.
The polynomials  $\Phi^i_{jl;mk}$ not involving the index $n$ describe
the base space $B_{n-1}$, and  by the induction hypothesis $\frac{n(n-1)(n-5)}6$  of them are linearly independent (the base case is $n=5$, where no such polynomials exist).
This gives   $\frac{n(n-3)}2+
\frac{n(n-1)(n-5)}6= \frac{n(n+1)(n-4)}6$ linear independent quadratic
equations for $B_n$. 
We remark that this procedure does not give polynomials of the form  $\Phi^n_{jl;km}$,
but equation \eqref{phi-eqn} shows that such polynomials are linear combinations
of the other ones. Therefore we obtain all equations of  Theorem \ref{versal-def-thm}
for $B_n$.
\end{proof}

\end{document}
